\subsectionaddtolist{\lr{Ubiquitous Language}}

زبان مشترک یا \lr{Ubiquitous Language} یکی از ارکان اصلی و حیاتی در رویکرد \lr{DDD} است. این مفهوم به یک واژه‌نامه و زبان دقیق اشاره دارد که توسط تمام اعضای تیم، شامل متخصصان کسب‌وکار (\lr{Domain Experts}) و توسعه‌دهندگان نرم‌افزار، به صورت توافقی ساخته و استفاده می‌شود. هدف اصلی این است که شکاف ارتباطی و سوءتفاهم‌های ناشی از ترجمه زبان تجاری به زبان فنی از بین برود. در این رویکرد، اصطلاحاتی که در جلسات و مکالمات روزمره تیم به کار می‌روند، باید دقیقا همان نام‌هایی باشند که در طراحی مدل‌ها، کلاس‌ها، متدها و حتی متغیرها در سورس‌کد استفاده می‌شوند. این همگام‌سازی تضمین می‌کند که نرم‌افزار دقیقا بازتاب‌دهنده نیازها و مفاهیم واقعی کسب‌وکار باشد.

\subsectionaddtolist{\lr{Domain Service , Application Service}}

در معماری \lr{DDD} برای جلوگیری از در هم آمیختگی وظایف، سرویس‌ها به دو لایه مجزا با اهداف متفاوت تقسیم می‌شوند.

سرویس دامنه (\lr{Domain Service}) دربرگیرنده آن دسته از قوانین و منطق‌های تجاری (\lr{Business Logic}) است که به طور ذاتی به یک موجودیت (\lr{Entity}) یا شیء مقداری (\lr{Value Object}) خاص تعلق ندارند. این سرویس‌ها معمولا عملیاتی را مدل می‌کنند که نیازمند تعامل بین چندین شیء مختلف در دامنه است. سرویس‌های دامنه کاملا درگیر هسته کسب‌وکار هستند و هیچ‌گونه حالت (\lr{State}) پایداری را در خود نگه نمی‌دارند.

در مقابل، سرویس کاربردی (\lr{Application Service}) فاقد هرگونه منطق تجاری است و صرفا نقش یک هماهنگ‌کننده (\lr{Orchestrator}) را ایفا می‌کند. وظیفه اصلی این سرویس، دریافت درخواست‌ها از لایه‌های بیرونی نرم‌افزار، مدیریت تراکنش‌ها، بررسی دسترسی‌ها و سپس فراخوانی اشیاء و سرویس‌های لایه دامنه برای انجام کار اصلی است. به زبان ساده، \lr{Use Cases} سیستم را پیاده‌سازی می‌کند تا دامنه بتواند بدون درگیری با مفاهیم زیرساختی، فقط روی منطق خالص خود تمرکز کند.

\subsectionaddtolist{\lr{Repository}}

الگوی مخزن یا \lr{Repository} ابزاری برای جداسازی منطق خالص دامنه از جزئیات فنی مربوط به ذخیره‌سازی و بازیابی داده‌ها (مانند دیتابیس‌ها) است. در دیدگاه \lr{DDD}، یک مخزن شبیه به یک مجموعه داده درون‌حافظه‌ای (\lr{In-Memory Collection}) عمل می‌کند که صرفا برای ذخیره و بازیابی ریشه‌های تجمیع (\lr{Aggregate Roots}) طراحی شده است. مدل دامنه نباید هیچ‌گاه درگیر نحوه نوشتن کوئری‌های \lr{SQL} یا ارتباط با \lr{ORM}ها شود؛ بلکه تنها از طریق رابط‌ها (\lr{Interfaces}) با مخزن صحبت می‌کند. مخزن در پس‌زمینه، پیچیدگی‌های ارتباط با زیرساخت داده را پنهان می‌کند و به توسعه‌دهندگان اجازه می‌دهد تا هسته نرم‌افزار را کاملا مستقل و ایزوله از تکنولوژی‌های دیتابیس نگه دارند.